% ============================================================
% Paper deep reading: Wan-Image
% ============================================================

\section{论文精读：Wan-Image}

\begin{conceptbox}
\textbf{论文：} \href{https://arxiv.org/abs/2604.19858}{Wan-Image: Pushing the Boundaries of Generative Visual Intelligence}\\
\textbf{作者/机构：} Wan Team, Alibaba Group\\
\textbf{时间：} arXiv v2, 2026-04-23；报告日期：2026-04-27；阅读日期：2026-04-29\\
\textbf{方向标签：} Image Generation、Image Editing、MLLM Planner、DiT Visualizer、4-channel VAE、Cascade RL、4K Refiner\\
\textbf{模型卡：} Wan2.7-Image；\href{https://create.wan.video/generate/image/generate?model=wan2.7-pro}{create.wan.video}
\end{conceptbox}

\subsection{0. 一句话结论}

Wan-Image 不是一篇“提出单个新模块”的论文，而是一份完整的工业级图像生成系统报告。它的目标不是证明某个纯架构假设，而是把图像模型从 casual image synthesizer 推向 professional productivity tool：复杂文字、身份保持、多图编辑、连续图组、透明 alpha、极端长宽比和 4K 输出都被当成一等能力来设计。

\begin{summarybox}
\textbf{我的判断：} Wan-Image 的核心价值在于系统工程路线：强 MLLM Planner 负责理解、路由和 prompt planning，DiT Visualizer 负责高保真生成，4-channel VAE 解决文字/透明图层/高频细节，Prompt Enhancer 和 Image Refiner 把用户意图与最终生产质量再拉一遍，Cascade RL 负责偏好和设计质量对齐。这篇不适合只读成“又一个 DiT”，而应该读成“专业视觉生产模型需要哪些基础设施”。
\end{summarybox}

\subsection{1. 它想解决什么问题}

论文开场的判断很直接：当前图像生成模型已经能做高审美图片，但还经常卡在生产工作流中的硬约束：

\begin{itemize}[leftmargin=1.5em]
  \item \textbf{复杂 typography：} 长文本、多行排版、海报和信息图中的文字仍容易错字、粘连或布局不稳。
  \item \textbf{严格 controllability：} 真实设计任务要求空间、逻辑、颜色、风格、身份和局部编辑都可控。
  \item \textbf{identity preservation：} 多参考主体、多图编辑、图组生成中，人物和物体身份不能漂移。
  \item \textbf{生产格式：} 电商、游戏、设计素材需要透明 alpha、极端长宽比和 2K--4K 质量，而不是单张 1K 美图。
\end{itemize}

因此 Wan-Image 的问题定义不是“如何提升文生图平均美感”，而是：

\begin{intubox}
如何把 MLLM 的语义理解、任务规划能力，与 DiT 的高质量像素合成能力，组合成一个面向专业生产的统一视觉生成系统？
\end{intubox}

\begin{figure}[htbp]
  \centering
  \includegraphics[width=0.82\linewidth]{wan_image/wan_image_fig1_eval_radar.pdf}
  \caption{Wan-Image 的综合人评雷达图。论文用 T2I/general editing 的 Win Score，以及 interactive editing/image series 的 Pass Rate 来展示能力覆盖。读这张图时要注意：Wan-Image 的卖点不是单项 benchmark，而是专业任务面上的覆盖，包括文字、身份、编辑、图组和设计场景。图源：arXiv source。}
\end{figure}

\subsection{2. 系统总览：Planner + Visualizer}

Wan-Image 的整体结构可以简化成：

\[
\begin{aligned}
\text{User Intent}
&\rightarrow \text{MLLM Planner}
\rightarrow \text{Prompt / Plan}\\
&\rightarrow \text{DiT Visualizer}
\rightarrow \text{Image Refiner}
\end{aligned}
\]

其中 Planner 与 Visualizer 不是简单串两个黑盒。论文称其采用 unified Transformer architecture：理解分支是 decoder-only Transformer，生成分支是 Rectified Flow 下训练的 DiT；二者在统一框架里对齐表示，并在训练上采用 decoupled strategy：先训练理解分支，再用理解分支初始化生成分支的 attention weights，新引入模块随机初始化。

\begin{figure}[htbp]
  \centering
  \includegraphics[width=\linewidth]{wan_image/wan_image_fig3_framework.pdf}
  \caption{Wan-Image 架构。重点看三层耦合：MLLM Planner 负责理解、推理和任务路由；DiT Visualizer 负责生成；VAE token、ViT token 和 text token 通过结构化 attention 组织。它不是 Tuna-2 那类 encoder-free 设计，而是更偏工业系统的 MLLM + ViT + VAE + DiT 组合。图源：arXiv source。}
\end{figure}

\subsubsection{2.1 Attention 与 token 组织}

论文里有几个关键设计：

\begin{itemize}[leftmargin=1.5em]
  \item \textbf{文本、ViT token、VAE token 之间用结构化 attention mask：} 文本侧保持 auto-regressive 生成，视觉 token 内部使用 bidirectional attention。
  \item \textbf{图组生成时，多个图片的视觉 token 之间也启用 bidirectional attention：} 这服务于 image series 的身份一致性和风格连贯性。
  \item \textbf{编辑任务里输入图像的 clean VAE tokens 作为条件输入：} clean tokens 的 timestep 显式设为 0，类似 teacher-forcing 的确定性条件；noised VAE tokens 则走生成分支。
  \item \textbf{位置编码分流：} 理解分支使用 MRoPE，生成分支使用 3D-RoPE，把 \(T,H,W\) 三个维度分开编码，并在 text/image segments 之间加固定 temporal offset 作为语义边界。
\end{itemize}

这里的重点是：Wan-Image 并不追求“把所有东西压成一个最纯的 token stream”。它更像把每种生产能力所需的 inductive bias 都保留下来，再通过统一 Transformer 和训练流程让它们协同。

\subsection{3. 4-channel VAE：为什么它还需要 VAE}

这篇和 Tuna-2 的立场几乎相反。Tuna-2 问的是“能否去掉 VAE 和 vision encoder”；Wan-Image 则认为，在生产级图像系统里，VAE 的质量直接决定生成上限，尤其是：

\begin{itemize}[leftmargin=1.5em]
  \item 小字号文字、密集排版、文档区域；
  \item 复杂纹理、细边界和高频局部结构；
  \item 透明 PNG、半透明边缘、图层素材和电商资产。
\end{itemize}

Wan-Image 提出 high-fidelity four-channel VAE，同时建模 RGB 和 alpha。结构上先做 \(2\times2\) patchify，再经过 \(8\times8\) spatial downsampling，总压缩率是 \(16\times16\)，latent dim 为 48。训练上采用三阶段：

\begin{enumerate}[leftmargin=1.5em]
  \item \(256\times256\) 低分辨率训练，先稳定透明内容和基本结构。
  \item mixed-resolution 训练，增强多尺度泛化。
  \item 加入 4-channel GAN discriminator，提升锐度、透明边界和整体一致性。
\end{enumerate}

此外它还加入 semantic distillation，让 latent 不只保重建像素，也携带更强语义信息。论文报告的重建结果：

\begin{center}
\small
\begin{tabular}{p{0.36\linewidth} p{0.20\linewidth} p{0.16\linewidth} p{0.16\linewidth}}
\toprule
\textbf{设置} & \textbf{PSNR} & \textbf{SSIM} & \textbf{LPIPS}\\
\midrule
RGB, 512p, 5K images & 35.429 & 0.958 & 0.010\\
RGBA, 1080p, 500 images & 38.274 & 0.970 & 0.018\\
\bottomrule
\end{tabular}
\end{center}

\begin{warnbox}
\textbf{和 Tuna-2 的区别：} Tuna-2 去掉 VAE 是为了验证 pixel-space unified modeling 的可行性；Wan-Image 保留并强化 VAE，是因为专业生产要可控、高清、可编辑、可透明合成。二者回答的问题不同，不能简单说谁路线更先进。
\end{warnbox}

\subsubsection{补充：VAE 是否应该保留}

如果把 Wan-Image 和近期一些去 VAE 的路线放在一起看，例如 Tuna-2、SenseNova 这类更偏 encoder-free / pixel-space 的方向，我的判断不是“VAE 必须保留”或“VAE 已经过时”，而是：\textbf{这取决于你到底想优化什么。}

\begin{itemize}[leftmargin=1.5em]
  \item \textbf{更倾向保留 VAE 的场景：} 目标里明确包含小字号文字、密集排版、透明 alpha、局部编辑、身份保持、素材合成、2K--4K 交付，或者你更在意商业工作流里的稳定边界。这类任务本质上都在惩罚“高频细节丢失”和“重建不可控”，而一个足够强的专用 VAE 可以把压缩、重建和语义先验先固定住，降低后端生成器需要同时解决的难题数量。
  \item \textbf{更倾向去掉 VAE 的场景：} 你更在意统一建模纯度、端到端 token 表示、理解与生成共享同一空间、避免 latent bottleneck，以及继续把 scaling law 往 pixel-space 推进。在这种目标下，VAE 本身反而可能变成中间瓶颈：它的压缩误差、频率偏置和重建上限，会提前决定系统能达到的视觉边界。
\end{itemize}

Wan-Image 代表的立场可以概括成一句话：\textbf{它愿意接受 latent bottleneck，换取更强的重建确定性和工程可控性。} 这对专业生产系统是合理的，因为它服务的不是“证明 unified modeling 能否完全去中介”，而是“让文字、透明图层、复杂编辑和高分辨率输出尽量稳定落地”。相对地，Tuna-2 / SenseNova 一类路线更像是在赌另一件事：随着模型规模、算力和 pixel-space 训练方法继续推进，VAE 这个中间压缩器本身，也许才是限制统一能力和端到端扩展性的地方。

\begin{intubox}
所以我的结论不是“VAE 一定要保留”，而是：如果目标是 \textbf{生产交付稳定性}，我会先保留强 VAE；如果目标是 \textbf{统一建模与长期 scaling 潜力}，去 VAE 是值得下注的路线，但要准备为文字、alpha、超高分辨率和编辑稳定性补更多额外工程。
\end{intubox}

\subsection{4. 数据引擎：真正的护城河}

Wan-Image 把数据系统写得很重。数据覆盖理解和生成两大部分。

\subsubsection{4.1 理解数据}

\begin{itemize}[leftmargin=1.5em]
  \item \textbf{General understanding data：} 文本和图文 QA，过滤、平衡、用 LLM/MLLM 重写低质量答案，并构造带 \texttt{<think>} 的 reasoning supervision。
  \item \textbf{Text-proxy data：} 为 interleaved generation 构造密集视觉 prompt。每个要生成的图片位置用 \texttt{<BOI>} 标记，视觉描述放在 \texttt{<imagine>} 与 \texttt{</imagine>} 内。
  \item \textbf{User-aligned query set：} 用 MLLM 按任务类型、难度、是否需要数值/逻辑推理、是否纯 VQA 等属性标注真实用户 query，再按规则采样，提升实用性。
\end{itemize}

\subsubsection{4.2 生成数据}

生成任务覆盖：

\begin{itemize}[leftmargin=1.5em]
  \item T2I：支持长文本渲染和最高 1:8 的极端长宽比。
  \item I2I：编辑、参考生成、多图条件，最多支持 9 张输入图。
  \item T2S / TI2S：一次生成最多 12 张图，强调视觉一致性和逻辑进展。
\end{itemize}

数据获取不是普通 web scrape，而是 multi-modal retrieval + annotation + cleaning 的闭环：支持 image-to-image、多图检索、text-to-image、图文混合和 batch retrieval，并用 cluster-based diversity re-ranking 控制多样性。

\begin{figure}[htbp]
  \centering
  \includegraphics[width=0.88\linewidth]{wan_image/wan_image_fig2_data_taxonomy.pdf}
  \caption{Wan-Image 的 structured data taxonomy。它先把图片路由到 photorealistic、non-photorealistic、text-centric、charts、multi-image compositions 五类，再按类别挂载不同属性维度。这个设计服务于两个目标：减少数据分布偏斜，并让模型学会专业任务需要的细粒度约束。图源：arXiv source。}
\end{figure}

\subsubsection{4.3 Annotation 与过滤}

论文的数据过滤 operator 包括五类：image feature extraction、aesthetic quality、AI-generated content detection、low-level information、overall image quality。具体例子包括：

\begin{itemize}[leftmargin=1.5em]
  \item compression artifact ratio：用理论未压缩大小和实际文件大小的比例判断过压缩。
  \item edge pixel variance：过滤大面积纯色背景和有效信息低的图片。
  \item JPEG bits-per-pixel：作为结构复杂度 proxy，过滤低信息密度样本。
  \item artificiality score：过滤明显 AI 生成或不真实图。
  \item watermark / greasy texture detection：提升数据可用性和真实感。
\end{itemize}

caption 也分三层：

\begin{itemize}[leftmargin=1.5em]
  \item raw captions：保留网页 alt-text、tag、标题等噪声弱监督。
  \item natural language captions：用 MLLM 生成多粒度自然语言描述。
  \item structured JSON captions：按类别填 25 个左右属性维度，再由 MLLM 重写成细节丰富的自然语言。
\end{itemize}

\subsection{5. Prompt Enhancer 与 Image Refiner}

\subsubsection{5.1 Prompt Enhancer}

Prompt Enhancer 是 Wan-Image 很工程化但很关键的一层。它不是单纯扩写 prompt，而是把用户意图翻译成更适合生成模型执行的设计 brief。

PE 有两个版本：

\begin{itemize}[leftmargin=1.5em]
  \item \textbf{Non-CoT：} 基于 Qwen3-VL-2B，输出约 400--600 tokens，面向低延迟，推理约 1 秒。
  \item \textbf{CoT：} 基于 Qwen3-VL-30B-A3B，输出约 1500--2000 tokens，显式做设计推理、元素扩展和空间布局规划。
\end{itemize}

它覆盖 T2I、I2I、T2S、TI2S。对 reference-guided I2I/TI2S，PE 输出 common 和 difference 两个字段：common 描述应保持不变的主体、身份或风格，difference 描述要修改的部分。这个拆分很重要，因为真实编辑任务的难点往往不是“改什么”，而是“什么绝对不能改”。

PE 的训练是 SFT + RL，RL 阶段使用 GRPO，reward 评估 intent fidelity、descriptive richness、linguistic accuracy。论文还提到 prompt rewriting 不是所有任务都强制开启，而是对文字渲染、海报、信息图、UI、storyboard、技术图解、逻辑推理等复杂类别选择性启用。

\subsubsection{5.2 Image Refiner}

Image Refiner 用在 DiT 输出之后，目标是把输出推到 2K--4K，并修复局部缺陷。它有两种模式：

\begin{itemize}[leftmargin=1.5em]
  \item \textbf{repair mode：} 输入生成图和详细 prompt，修复错字、噪声、局部 artifact。
  \item \textbf{high-fidelity mode：} 用在 identity preservation，尽量减少结构和身份漂移。
\end{itemize}

训练时用 blur、noise、lossy compression，以及 diffusion forward-noising/denoising 构造退化输入；之后再用 distillation 和 RL 优化效率与质量。

\subsection{6. 训练流程}

Wan-Image 的训练分为理解、生成、RL、distillation 四段。

\subsubsection{6.1 理解训练}

理解部分从 Qwen-VL base model 初始化，三阶段：

\begin{itemize}[leftmargin=1.5em]
  \item \textbf{Understanding CT：} 46K steps，0.1T tokens，混合图文理解和纯文本数据，保留语言推理能力。
  \item \textbf{Understanding \& Multi-task SFT：} 10K steps，0.033T tokens，用 understanding data 和 text-proxy data 联合训练，让模型能把用户输入转成生成指导。
  \item \textbf{On-policy multi-teacher distillation：} 0.4K steps，面向真实用户 query，用不同 teacher 迁移不同输入类别的处理能力。
\end{itemize}

\subsubsection{6.2 生成训练}

\begin{center}
\footnotesize
\setlength{\tabcolsep}{3pt}
\begin{tabular}{p{0.13\linewidth} p{0.13\linewidth} p{0.15\linewidth} p{0.20\linewidth} p{0.24\linewidth}}
\toprule
\textbf{阶段} & \textbf{Steps} & \textbf{Tokens} & \textbf{Resolution} & \textbf{Ratio}\\
\midrule
PT & 713K & 13.27T & 192 / 320 / 640 & 7:3, T2I/I2I\\
CT & 223K & 8.85T & 512--2048 & 7:2:0.5:0.5\\
SFT & 13K & 0.62T & 512--2048 & curated mix\\
\bottomrule
\end{tabular}
\end{center}

读法很清楚：PT 先靠大规模低中分辨率数据学通用生成；CT 把分辨率拉到 512--2048，并加入 T2S/TI2S，比例为 T2I:I2I:T2S:TI2S = 7:2:0.5:0.5；SFT 用更少但更干净的数据提升审美、细节和指令遵循。

\begin{figure}[htbp]
  \centering
  \includegraphics[width=0.92\linewidth]{wan_image/wan_image_fig5_training_stages.pdf}
  \caption{PT、CT、SFT 阶段输出对比。论文想表达的是训练阶段的职责分工：PT 学基本生成，CT 补高分辨率和复杂条件，SFT 强化美学和局部细节。图源：arXiv source。}
\end{figure}

\subsubsection{6.3 Cascade RL}

RL 阶段是这篇和你后训练专区连接最强的部分，而且它的 framing 很值得注意：Wan-Image 并不把 RL 写成“最后上一遍同一种偏好优化”，而是把它做成一个 \textbf{按 domain 组织、按 reward 形态选工具、按风险逐层推进} 的级联系统。这里的关键词不是某一个算法名，而是 \textbf{Cascade}。

先说论文明写的两条基础原则。第一，Wan-Image 构造 \textbf{domain-wise preference data}，说明偏好数据不是混成一个大池子，而是围绕不同生成场景分别采集和筛选。第二，\textbf{RL prompts 与 SFT prompts 互斥}，这是很重要的反过拟合设计：如果 RL 继续在 SFT 看过的 prompt 上做，很容易把 reward 学成“记住训练分布里什么图更像标注者喜欢”，而不是学到可迁移的偏好对齐能力。

再看图 4，论文把 RL 前面的 data/reward pipeline 单独画了出来：\textbf{Data Collection}、\textbf{Data Pair}、\textbf{Aesthetic Expert}、\textbf{Data Retrieval}、\textbf{Quality Monitoring}、\textbf{Reward Benchmark}、\textbf{Reward Model}。这其实透露了 Wan-Image 的系统思路：\textbf{先做偏好数据工程，再做 reward 工程，最后才是 RL 工程}。也就是说，它并不把 reward model 当成一个“训完就用”的静态打分器，而是持续用更广的 benchmark 和人工质检去检查 reward 是否偏离真实偏好。

从图里的维度名还能看出 reward 不是单一 aesthetic score，而是至少覆盖：\textbf{Realism、Text、Harmony、Structure、Aesthetics、Consistency}。这和前面模型的产品目标是对齐的：

\begin{itemize}[leftmargin=1.5em]
  \item \textbf{Realism：} 针对写实感、材质和整体自然度；
  \item \textbf{Text：} 针对文字可读性、拼写与排版；
  \item \textbf{Structure：} 针对空间关系、布局和物体结构正确性；
  \item \textbf{Aesthetics：} 针对构图、色彩和设计感；
  \item \textbf{Consistency / Harmony：} 更偏编辑、多图和组合场景里的身份保持、风格一致与局部协调。
\end{itemize}

这个设计的直接含义是：\textbf{Wan-Image 的 RL 目标不是把所有样本都往“更漂亮”一个方向推，而是按任务暴露出的短板做多维对齐。} 这也解释了为什么它最后不是单一 RL，而是 DPO、ReFL、DenseGRPO、ReFMA 的组合。

Cascade RL 组合了：

\begin{itemize}[leftmargin=1.5em]
  \item DPO；
  \item ReFL；
  \item DenseGRPO；
  \item in-house ReFMA，即 reinforcement learning with flow matching anchors。
\end{itemize}

我的理解是，这四类方法并不是互相替代，而是分别覆盖不同风险/成本/信用分配条件下的优化位置。

\textbf{1. DPO：先用离线 pairwise preference 把方向拉正。}

DPO 的优点是稳定、便宜、工程风险低。给定同一 prompt 下的 winner/loser 对，模型直接学“让 winner 相对 loser 更可能”，同时参考模型提供一个 KL-like anchor，避免分布漂太远。放到 Wan-Image 这里，DPO 很适合作为级联里的低风险起点：\textbf{先把明显的人类偏好方向学进去，例如更好的排版、更合理的构图、更顺眼的整体风格，而不急着上高成本在线探索。}

但 DPO 的边界也很明确：它主要吃现有偏好对，\textbf{不会自己探索新轨迹}，也不能细致回答“到底是哪几个 denoising steps 把图像变好了”。所以它更像是偏好对齐的第一层，不太像最终解法。

\textbf{2. ReFL：当 reward 可微时，直接把 reward 梯度压回生成过程。}

从 Wan-Image 把 ReFL 放进同一工具箱来看，它显然想利用一类和 policy-ratio 不同的优化信号：\textbf{如果 reward model 可微，就可以不只靠采样后比较 winner/loser，也可以直接沿 reward gradient 优化生成轨迹。} 这类方法的直觉通常是：对那些 reward 能稳定评估、而且修改更像局部连续修正的目标，直接梯度会比纯采样式 RL 更高效。

结合你前面后训练主线里的经验，可以把 ReFL 粗略记成“更偏 direct reward optimization 的路线”。它通常更容易先改善后期可见的质量问题，比如局部细节、审美微调、文字边缘和表面纹理。但它的前提也更苛刻：\textbf{reward 要可微，训练稳定性也更依赖 reward 本身的质量。}

\textbf{3. DenseGRPO：解决 flow/diffusion 里的时间步信用分配。}

这部分最值得和你已经读过的 DenseGRPO 对照。标准 Flow-GRPO/GRPO 的一个硬伤是：最终图像只有一个 terminal reward，但整条生成轨迹有很多 denoising steps。如果把同一个终局 reward 平均分给所有 timestep，就等于默认“最终图像好，说明每一步都好”，这通常不成立。

DenseGRPO 的价值正是在这里。它不满足于 trajectory-level reward，而是试图估计每一步到底带来了多少 reward gain。具体做法是：从中间 latent 通过 ODE readout 还原出可打分的 clean prediction，再看相邻 step 的 reward 提升量。这样得到的是 step-wise dense reward，而不是把同一个 terminal reward 广播到全轨迹。

放到 Wan-Image 的语境里，DenseGRPO 很适合那些\textbf{明显依赖生成过程信用分配}的目标，例如结构稳定性、文字渲染和多对象空间关系。因为这些问题往往不是“最后一张图总体不好看”这么简单，而是某些关键步把 layout、shape 或字符骨架画错了。如果 credit assignment 太粗，模型只知道“这张图分低”，却不知道应该改前期布局还是后期 refinement。

当然，DenseGRPO 的代价也很实在：中间步读出和打分很贵，工程上明显重于 DPO。Wan-Image 愿意把它放进 Cascade RL，本身就说明他们认为单靠离线偏好或整图 reward 已经不够解决生产场景里的细粒度问题。

\textbf{4. ReFMA：用 flow-matching anchor 稳住 reward-driven optimization。}

论文对 ReFMA 的公开细节没有 DenseGRPO 那么充分，但从名字和图示可以读出两个关键信号。第一，它是 \textbf{flow matching anchors}，说明它不是完全放飞的 reward maximization，而是保留一个和原始 flow trajectory/teacher policy 相关的锚点。第二，现有描述强调 \textbf{在不同 noise levels 之间平衡优化目标}，并使用 \textbf{reward model 的梯度信号} 指导优化。

这意味着 ReFMA 很可能是在解决一个非常实际的问题：\textbf{直接用 reward 去推 flow model，容易在某些噪声区间过度优化，导致分布漂移、结构塌陷或者局部 reward hacking；anchor 的作用就是把模型拴回原始可用分布附近。} 如果这么理解，ReFMA 扮演的是“高强度 reward optimization 的安全带”。

我会把它看成 Wan-Image 在工业落地里的关键补件：\textbf{既想吃到 reward-gradient 带来的强对齐收益，又不愿意让模型因为几类 reward 被拉坏基础生成分布。} 这和纯研究论文常见的“只报 reward 提升”不同，Wan-Image 更在乎长期可用性和跨任务稳定性。

\textbf{把四者连起来看，Cascade RL 背后的逻辑就清楚了。}

如果只用 DPO，优点是稳，但缺少在线探索和 step-level credit；如果只用 DenseGRPO，credit 更准，但成本高；如果只用 ReFL 一类 direct-gradient 方法，优化效率高，但受 reward 可微性和稳定性约束；如果只有强 reward gradient 而没有 anchor，又容易过冲。Wan-Image 把它们级联，本质上是在做下面这件事：

\begin{itemize}[leftmargin=1.5em]
  \item 先用更稳的离线偏好方法把大方向拉正；
  \item 再用更强的 reward-based 方法补足在线探索和细粒度 credit；
  \item 最后用 anchor/control 机制把 reward optimization 限制在可用分布附近。
\end{itemize}

所以这部分最值得记住的，不是“Wan-Image 也用了 DPO/DenseGRPO”，而是：

\begin{intubox}
对于专业图像系统，RL 不是一个统一算法，而是一套要和数据分区、reward 维度、生成轨迹结构、稳定性约束一起设计的后训练工具链。
\end{intubox}

论文报告 Cascade RL 相比 SFT baseline 在不同场景下 Win Score 提升约 15\%--50\%。结合图 4 的 reward 维度，这个提升不应理解成单纯 aesthetic bump，而更像是对 \textbf{photorealism、structural correctness、text rendering、aesthetic quality，以及 I2I consistency/harmony} 这些生产敏感指标的系统性补强。

\begin{figure}[htbp]
  \centering
  \includegraphics[width=0.92\linewidth]{wan_image/wan_image_fig4_rl_framework.pdf}
  \caption{Cascade RL 框架和相对 SFT baseline 的提升。图里除了 DPO、ReFL、DenseGRPO、ReFMA 的级联关系，还明确画出了 reward system 的前置基础设施：data pair、aesthetic expert、data retrieval、quality monitoring、reward benchmark、reward model，以及 realism/text/harmony/structure/aesthetics/consistency 六个对齐维度。图源：arXiv source。}
\end{figure}

\subsubsection{6.4 Distillation}

蒸馏在 Wan-Image 里不是泛指的“模型压缩”，而是非常具体的 \textbf{inference step distillation}：目标不是再造一个更小的基座模型，而是在尽量不破坏生成能力的前提下，把原始采样流程压缩到更少 step，从而满足真实交互应用对低延迟的要求。论文特别指出，尽管 base model 的参数规模已经相对克制，但其原始的 iterative generation scheme 再叠加 classifier-free guidance (CFG)，仍然会带来显著计算负担；部署时真正昂贵的，不只是参数量本身，而是一次生成要做多少次 denoising/function evaluation，也就是 NFE。

这里最重要的理解是：\textbf{Wan-Image 要解决的不是一般意义上的“少步采样会降质”，而是少步采样会优先破坏那些最有产品价值的能力。} 论文明确点出，已有 step distillation 方法虽然可以提升速度，但往往会在 \textbf{high-frequency facial details} 和 \textbf{text rendering capability} 上出现明显退化。这两个点之所以值得作者单独拿出来说，是因为它们恰好对应 Wan-Image 想重点建立优势的使用场景：真实人像、广告图文、海报排版、品牌素材，以及需要多轮编辑的商业工作流。也就是说，Wan-Image 的蒸馏目标不是“把图粗略生成出来”就够了，而是要求在提速之后，仍尽量保住脸部微细节、文字可读性、局部纹理锐度，以及整体结构稳定性。

为了解这个矛盾，论文没有采用单一蒸馏技巧，而是设计了一个 \textbf{multi-stage progressive distillation} 框架，把两类互补思路组合起来：trajectory distillation 和 distribution matching distillation。虽然这一节没有公开每个 stage 的精确步数、loss 权重和 scheduler 细节，但从作者对目标的描述可以比较清楚地看出，两者扮演的角色并不一样。

\begin{itemize}[leftmargin=1.5em]
  \item \textbf{Trajectory distillation：} 让 student 用更少的采样步去逼近 teacher 的 denoising trajectory。按我的理解，它更偏向“过程对齐”，核心作用是保住全局布局、主体关系、结构演化，以及多步去噪过程中逐步收敛的动态。对于高容量 teacher 来说，这一步尤其重要，因为 teacher 的去噪轨迹通常比小模型更复杂，如果直接从多步 teacher 硬压到极少步 student，学生模型很容易只学到一个粗糙终态，结果就是构图漂移、局部关系不稳，或者 prompt 条件在中间过程里被过早丢失。
  \item \textbf{Distribution matching distillation：} 进一步把 student 的输出分布往 teacher 的高质量分布上拉齐。它更像“结果对齐”，关注的不是每一步都长得像 teacher，而是最后生成出来的样本在统计意义上仍然落在 teacher 擅长的那类结果附近。结合论文强调的退化点，这一步可以理解为对 trajectory imitation 的补充：前者负责守住结构与轨迹，后者负责补 realism、纹理锐度和文本 fidelity，尽量避免少步近似之后出现的过平滑、字形塌陷、局部细节发糊等问题。
\end{itemize}

为什么还要做成 \textbf{progressive} 的多阶段蒸馏？一个合理解释是，作者并不想让 student 一次性从“高 NFE teacher”直接跳到“极少步部署模型”，而是把压缩过程做成一种 curriculum。这样每个阶段只承担一部分步数压缩和误差吸收，训练更容易稳定，也更容易在速度和保真之间找到可接受的折中。论文虽然没有展开完整训练配方，但从问题设定上看，这种 staged compression 本质上是在规避一个老问题：单次激进压步通常能换来更快采样，却很容易把 fidelity 尤其是细节 fidelity 一起压坏。

所以这部分最值得记住的，不是“Wan-Image 也做了蒸馏”，而是：

\begin{intubox}
Wan-Image 把蒸馏当成“保持能力边界的采样压缩”，而不是单纯追求更少 step。真正要守住的是 \textbf{全局结构、高频细节和文本对齐} 这三件事。
\end{intubox}

把它放回整篇论文的系统视角里看，这一步其实是前面所有能力模块的产品化收口：VAE 负责高压缩下的视觉重建，thinking/instruct 负责复杂理解与编辑控制，Cascade RL 负责把生成偏好往高质量分布推，而 distillation 则负责把这些能力压进真实可部署的延迟预算里。论文给出的结论是，经过这一套 multi-stage progressive distillation 后，推理流程获得了超过 \(8\times\) 的 speedup，同时仍能较好保持 generation fidelity。这里的“保真”更应理解为系统级保真，而不是所有子任务严格无损；但从作者特意强调脸部高频细节、文本渲染，以及 Figure 19 中 teacher vs. distilled model 的对比可以看出，他们真正担心的，就是那些最容易在少步采样里首先塌掉、却又最影响商业可用性的能力。
\subsection{7. 结果怎么读}

论文报告的核心结果包括：

\begin{itemize}[leftmargin=1.5em]
  \item \textbf{人评：} 整体超过 Seedream 5.0 Lite 和 GPT Image 1.5，在困难任务上接近 Nano Banana Pro。
  \item \textbf{交互编辑和图组生成：} pass rate 大约 80\%，强调真实工作流可用性。
  \item \textbf{理解能力：} Ours Thinking 在 7 个 image understanding benchmark 的平均分为 76.3，高于很多 unified models；text understanding 平均也高于其 instruct mode。
  \item \textbf{VAE：} RGB/RGBA 重建指标优于或接近强 baseline，尤其强调 \(16\times16\) 高压缩下的细节与 alpha。
  \item \textbf{RL：} Cascade RL 对 SFT baseline 有 15\%--50\% Win Score 提升。
\end{itemize}

但要谨慎：这篇没有给出完全可复现的开放训练配方，也没有把所有模块做严格公开 ablation。很多结论来自系统级人评和质性展示，适合判断产品能力和工程方向，不适合直接当作单一算法的严格因果证明。

\subsection{8. 与 Tuna-2 / BAGEL 的关系}

\begin{center}
\footnotesize
\setlength{\tabcolsep}{3pt}
\begin{tabular}{p{0.15\linewidth} p{0.25\linewidth} p{0.23\linewidth} p{0.25\linewidth}}
\toprule
\textbf{模型} & \textbf{视觉表示} & \textbf{生成路径} & \textbf{核心取舍}\\
\midrule
Tuna-2 & pixel embeddings；无 VAE/ViT & pixel-space flow & 架构假设更激进，证明 encoder-free scaling；生成侧未必最强。\\
BAGEL & SigLIP + VAE latent & Flow Matching + MoT & 理解/生成分路，强调 native unified。\\
Wan-Image & ViT + VAE + MLLM Planner & DiT + Refiner & 工业系统路线，保留并强化视觉先验，服务专业生产约束。\\
\bottomrule
\end{tabular}
\end{center}

Wan-Image 说明一个很现实的点：即使 encoder-free/pixel-space 是很有研究价值的方向，生产级图像系统仍可能选择强 VAE、强 ViT、强 MLLM、强 refiner 的组合。因为业务目标不是架构纯度，而是文字、身份、透明图层、编辑和 4K 交付的稳定性。

\subsection{9. 局限与风险}

\begin{itemize}[leftmargin=1.5em]
  \item \textbf{不可复现：} 数据规模、过滤器、reward model、teacher model 和私有标注流程都高度闭源。
  \item \textbf{系统归因困难：} 最终质量来自 PE、Planner、VAE、DiT、Refiner、RL、distillation 多模块叠加，很难知道单个模块贡献。
  \item \textbf{人评标准不透明：} Win Score 和 Pass Rate 有参考意义，但任务分布、标注一致性和失败样例覆盖仍需更细披露。
  \item \textbf{PE 可能改变用户意图：} prompt rewriting 提升美学和细节，但也可能把原始需求过度设计化。
  \item \textbf{RL 可能强化审美模式：} Cascade RL 会提升人类偏好分数，但也可能收缩风格多样性或偏向 reward model 的偏好。
  \item \textbf{4K 质量部分依赖 refiner：} 需要区分 base generation capability 和 post-refinement capability。
\end{itemize}

\subsection{10. 对我的可复用点}

\begin{itemize}[leftmargin=1.5em]
  \item \textbf{生产级指标先行：} 文字、身份、alpha、图组、局部编辑、极端比例，比普通 aesthetic prompt 更能暴露模型短板。
  \item \textbf{common/difference 编辑表示：} 对 reference editing 很实用，明确“保留什么”和“修改什么”。
  \item \textbf{数据 taxonomy + structured JSON caption：} 对图像生成训练比单纯自然语言 caption 更可控，尤其适合长尾任务。
  \item \textbf{RL 不是单算法：} 可以把 DPO、ReFL、DenseGRPO、flow-matching anchor 类方法按领域级联，而不是期待一个 reward objective 解决所有偏好。
  \item \textbf{VAE 是否保留取决于目标：} 如果目标是透明图层、小字、文档和 4K，专用 VAE 的工程价值很高；如果目标是追求 encoder-free 的统一建模与长期 scaling，可以继续关注 pixel-space 路线。
\end{itemize}

\subsection{11. 待查问题}

\begin{itemize}[leftmargin=1.5em]
  \item ReFMA 的具体 objective、anchor 设计和与 DenseGRPO 的衔接，论文没有完全展开，需要等技术细节或代码。
  \item unified Transformer 中 dedicated experts 的参数共享边界、attention 共享方式和推理时调度细节，需要进一步核对。
  \item Prompt Enhancer 的选择性触发规则是否会导致用户意图漂移，最好看更多实际案例。
  \item Refiner 对身份保持和文字修复的失败模式还不清楚，尤其是是否会在高分辨率增强时改动关键内容。
  \item 与 Nano Banana Pro、GPT Image 1.5、Seedream 的公平对比需要更透明的 prompt 集和人评协议。
\end{itemize}

\subsection{12. 后续阅读顺序}

\begin{itemize}[leftmargin=1.5em]
  \item \textbf{先对照 Wan2.2 / Wan video：} 看 Wan 系列的 VAE、DiT 和视频生成技术如何迁移到 image。
  \item \textbf{再读 DenseGRPO / Flow-GRPO / ReFL：} 补齐 Cascade RL 背后的偏好对齐工具。
  \item \textbf{对照 Tuna-2：} 前者看 encoder-free pixel scaling；后者看 VAE、ViT、MLLM 的生产系统栈。
  \item \textbf{对照 BAGEL：} 比较 native unified multimodal architecture 与 Wan-Image 的 planner-visualizer 工程路线。
\end{itemize}
